% generated by roam.py handbook export -- не редактировать, править заметки
% 20260922140145-неравенство_иенсена.org (body 8aec188530db)

\begin{theorem}[{Неравенство Йенсена}]\label[theorem]{note:53558BAA-700E-46BE-9EAB-6A2F65A5B4C0}
  Пусть $f$~--- \hyperref[note:1CA14F2C-637F-41DE-9E97-72075D19E97F]{выпуклая} на $E$, тогда
  \[
  \forall x_i \in E, 0 \le \lambda_i, \sum\limits_{i=1}^n \lambda_i = 1\quad f\left(\sum\limits_{i=1}^n \lambda_i x_i\right) \le
  \sum\limits_{i=1}^n \lambda_i f(x_i)
  \]
\end{theorem}

\begin{proof}
  Идея аналогична доказательству: \hyperref[note:42DFF6DC-CC9D-40AA-9D6F-218DCF613833]{Характеристика выпуклого множества через выпуклые комбинации}

  Доказывать будем по индукции.
  База: $n=1$,  тогда $\lambda_1 = 1 \Rightarrow f(\lambda_1 x_1) \le \lambda_1 f(x_1) \Leftrightarrow f(x_1) \le f(x_1)$. Очевидно;

  Давайте рассмотрим переход: $n \rightarrow n+1$ \\
  Имеем $\lambda_i \ge 0,\ \sum\limits_{i=1}^{n+1} \lambda_i = 1$. Если $\lambda_{n+1} = 1$, то мы опять вернемся к вырожденному случаю из базы,
  так как остальные $\lambda_i = 0, i < n+1$. Далее будем считать, что $\lambda_{n+1} \ne 1$.

  Пусть $s:= 1 - \lambda_{n+1}$, тогда
  \[
  \sum\limits_{i=1}^{n+1} \lambda_i x_i = \sum\limits_{i=1}^n \lambda_i x_i + \lambda_{n+1} x_{n+1} =
  s \cdot \frac{1}{s} \cdot \sum\limits_{i=1}^n \lambda_i x_i + \lambda_{n+1} x_{n+1}
  \]

  Рассмотрим $\frac{1}{s} \sum\limits_{i=1}^n \lambda_i x_i$. Это комбинация из $n$ точек, рассмотрим коэффициенты
  $\sum\limits_{i=1}^n \frac{\lambda_i}{s} = \frac{1-\lambda_{n+1}}{1-\lambda_{n+1}} = 1$, причем они все неотрицательные, так как $\lambda_i \ge 0$ и $s > 0$.

  Так как \hyperref[note:42DFF6DC-CC9D-40AA-9D6F-218DCF613833]{выпуклое множество содержит все свои выпуклые комбинации}, то $y := \sum\limits_{i=1}^n \frac{\lambda_i}{s} x_i \in E$.

  Итого, мы получили, что
  \[
  \sum\limits_{i=1}^{n+1} \lambda_i x_i = s y + \lambda_{n+1}x_{n+1} = (1-\lambda_{n+1}) y + \lambda_{n+1}x_{n+1}
  \]

  Воспользуемся определением \hyperref[note:1CA14F2C-637F-41DE-9E97-72075D19E97F]{выпуклости функции}.
  \[
  f( (1 - \lambda_{n+1})y + \lambda_{n+1} x_{n+1} ) \le (1-\lambda_{n+1}) f(y) + \lambda_{n+1} f(x_{n+1})
  \]
  Давайте теперь рассмотрим $f(y)$, это функция от \hyperref[note:8583F986-A6E8-4BBB-B2FB-BA0E9BBE58BF]{выпуклой комбинации} из $n$ точек.
  Значит мы можем применить индукционное предположение:
  \[
  f(y) = f \left(  \sum\limits_{i=1}^n \frac{\lambda_i}{s} x_i \right) \le \sum\limits_{i=1}^n \frac{\lambda_i}{s} f(x_i)
  \]

  Но тогда $(1-\lambda_{n+1}) f(y) = s f(y) \le \sum\limits_{i=1}^n \lambda_i f(x_i)$ \\
  Итого
  \[
  f \left( \sum\limits_{i=1}^{n+1} \lambda_i x_i  \right) = f(sy + \lambda_{n+1} x_{n+1}) \le sf(y) + \lambda_{n+1} f(x_{n+1}) \le \sum\limits_{i=1}^n \lambda_i f(x_i) + \lambda_{n+1} f(x_{n+1}) = \sum\limits_{i=1}^{n+1} \lambda_i f(x_i)
  \]
\end{proof}
